\documentclass[11pt]{article}

\usepackage[margin=0.75in]{geometry}
\usepackage[T1]{fontenc}
\usepackage[utf8]{inputenc}
\usepackage{lmodern}
\usepackage{microtype}
\usepackage{amsmath,amssymb}
\usepackage{booktabs}
\usepackage{tabularx}
\usepackage{array}
\usepackage{xcolor}
\usepackage{enumitem}
\usepackage{fancyhdr}
\usepackage{titlesec}
\usepackage{hyperref}
\usepackage{listings}

\definecolor{accent}{HTML}{1F4E79}
\definecolor{lightaccent}{HTML}{EAF2F8}
\definecolor{rulegray}{HTML}{C9D6E2}
\definecolor{textgray}{HTML}{333333}

\hypersetup{
  colorlinks=true,
  linkcolor=accent,
  urlcolor=accent,
  citecolor=accent,
  pdftitle={Point vs. Vertex: Formal Distinctions Across Graph Theory, Computational Geometry, and Meshes},
  pdfauthor={ChatGPT},
  pdfsubject={Mathematics and computer graphics terminology handout},
  pdfkeywords={point, vertex, node, graph theory, computational geometry, mesh, OpenGL, WebGL}
}

\pagestyle{fancy}
\fancyhf{}
\lhead{Point vs. Vertex}
\rhead{Terminology Handout}
\cfoot{\thepage}
\renewcommand{\headrulewidth}{0.4pt}
\renewcommand{\headrule}{\hbox to\headwidth{\color{rulegray}\leaders\hrule height \headrulewidth\hfill}}

\titleformat{\section}{\large\bfseries\color{accent}}{\thesection}{0.6em}{}
\titleformat{\subsection}{\normalsize\bfseries\color{accent}}{\thesubsection}{0.6em}{}
\setlist[itemize]{leftmargin=1.4em, topsep=3pt, itemsep=2pt}
\setlist[enumerate]{leftmargin=1.6em, topsep=3pt, itemsep=2pt}
\setlength{\parindent}{0pt}
\setlength{\parskip}{6pt}
\setlength{\headheight}{24pt}
\renewcommand{\arraystretch}{1.25}

\newcommand{\callout}[2]{%
  \vspace{4pt}%
  \noindent\colorbox{lightaccent}{%
    \parbox{0.965\linewidth}{%
      \textbf{\color{accent}#1}\par
      #2%
    }%
  }%
  \vspace{4pt}%
}

\newcolumntype{Y}{>{\raggedright\arraybackslash}X}
\newcolumntype{L}[1]{>{\raggedright\arraybackslash}p{#1}}

\begin{document}

\begin{center}
  {\LARGE\bfseries\color{accent} Point vs. Vertex}\\[3pt]
  {\large Formal Distinctions Across Graph Theory, Computational Geometry, and Meshes}\\[8pt]
  {\normalsize A concise handout for mathematical modeling, geometric algorithms, and graphics data structures}
\end{center}

\vspace{4pt}
\hrule height 0.6pt
\vspace{10pt}

\callout{Core thesis}{In casual communication, \emph{point}, \emph{vertex}, and sometimes \emph{node} are used loosely. Formally, the right term depends on whether you are talking about raw coordinate locations, abstract graph elements, or structured geometric/topological data.}

\section{Executive Summary}

A \textbf{point} is primarily a location in a coordinate space, such as
\[
  p=(x,y) \quad \text{or} \quad p=(x,y,z).
\]

A \textbf{vertex} is an element that plays a structural role in a larger object. Depending on context, that object may be a graph, polygon, polyhedron, triangulation, or mesh.

The practical rule is:
\[
  \text{Every geometric vertex has a point position, but not every point is a vertex.}
\]

That rule should be qualified for graph theory: a graph vertex is usually \emph{drawn} as a point, but the formal object is an abstract member of the graph's vertex set, not necessarily a coordinate-bearing geometric point.

\section{Graph Theory: A Vertex Is Abstract}

In graph theory, a graph is commonly written as
\[
  G=(V,E),
\]
where \(V\) is the set of vertices and \(E\) is the set of edges. In an undirected simple graph, an edge is usually modeled as an unordered pair of vertices. In a directed graph, an edge is commonly modeled as an ordered pair of vertices. This is the standard formal framing used in graph theory lecture notes and textbooks.\cite{stanford_graph_basics,nagoya_graph_notes,grinberg_graph_notes}

A graph vertex may be represented visually as a dot or point in a diagram, but this drawing convention does not make the vertex a geometric coordinate. Unless the graph is embedded into a geometric space, a vertex has no inherent \(x\), \(y\), or \(z\) location.

\subsection*{Precise wording}

\begin{itemize}
  \item Informal: ``The graph has five points connected by lines.''
  \item Better: ``The graph has five vertices connected by edges.''
  \item Most precise: ``The graph \(G=(V,E)\) has \(|V|=5\) vertices and edges specified by \(E\).''
\end{itemize}

\callout{Important distinction}{In graph drawings, vertices are often depicted as points. Formally, however, a vertex is an abstract graph element. Its visual location is part of a drawing or embedding, not part of the graph unless the model explicitly includes geometry.}

\section{Computational Geometry: Points Are Coordinates; Vertices Are Structural}

In computational geometry, a \textbf{point} is usually a coordinate-valued primitive. A set of points may be written as
\[
  P = \{p_1,p_2,\ldots,p_n\}, \qquad p_i \in \mathbb{R}^2 \text{ or } \mathbb{R}^3.
\]
Computational geometry studies algorithms over geometric objects such as points, line segments, polygons, triangulations, convex hulls, and Voronoi diagrams.\cite{mount_comp_geom,preparata_shamos_comp_geom}

A \textbf{vertex} is a point that has a role in a geometric or topological structure. Examples include:

\begin{itemize}
  \item a corner of a polygon;
  \item an endpoint of a line segment;
  \item a point incident to edges in a triangulation;
  \item a corner of a polyhedron;
  \item a boundary point selected by a convex hull algorithm.
\end{itemize}

For example, a point cloud may contain 1,000 raw coordinate samples. Those samples are just points until they are connected or selected into a structure. If a subset becomes the corners of a polygon, the vertices of a triangulation, or the boundary points of a convex hull, then those points now serve as vertices in that structure.

\section{Computer Graphics and Meshes: A Vertex Is Often a Record}

In computer graphics, a \textbf{vertex} is often more than a position. Graphics APIs describe vertex data through vertex attributes: position, normal vector, color, texture coordinates, tangents, or other shader inputs. WebGL's \texttt{vertexAttribPointer()} API, for example, specifies the layout of a generic vertex attribute in a bound vertex buffer.\cite{mdn_webgl_vertexattribpointer}

A typical mesh vertex record may be modeled as:
\[
  v_i = \{\text{position},\ \text{normal},\ \text{texcoord},\ \text{color},\ \text{tangent}\}.
\]

The \textbf{position} field is the point-like coordinate. The \textbf{vertex} is the structured record used by the rendering pipeline or mesh representation. OpenGL specifications also distinguish vertex-related data such as positions, normals, colors, and texture coordinates as part of graphics processing.\cite{opengl_spec}

\section{Decision Table}

\begin{tabularx}{\linewidth}{L{0.25\linewidth} Y Y}
\toprule
\textbf{Context} & \textbf{Use this term} & \textbf{Reason} \\
\midrule
Raw coordinate sample & Point & The object is only a location, such as \((x,y)\) or \((x,y,z)\). \\
Point cloud & Point & Samples usually have no required edges, faces, or incidence structure. \\
Graph theory & Vertex & The object is an abstract member of \(V\) in \(G=(V,E)\). \\
Network/system diagram & Node or vertex & ``Node'' is common in systems and networks; ``vertex'' is the graph-theoretic term. \\
Polygon or polyhedron & Vertex & The point is a corner or endpoint participating in a shape. \\
Triangulation or mesh topology & Vertex & The point participates in adjacency, incidence, edges, and faces. \\
Shader input / graphics buffer & Vertex attribute / vertex & A vertex may be a structured record containing position, normal, UV coordinates, and other attributes. \\
\bottomrule
\end{tabularx}

\section{Worked Examples}

\subsection*{Example 1: Raw coordinate data}

\begin{lstlisting}[basicstyle=\ttfamily\small, frame=single]
points = [(0, 0), (1, 3), (4, 2), (6, 5)]
\end{lstlisting}

These are best called \textbf{points}. There is no declared graph, polygon, triangulation, or mesh topology.

\subsection*{Example 2: Graph model}

\begin{lstlisting}[basicstyle=\ttfamily\small, frame=single]
V = {A, B, C, D}
E = {(A, B), (B, C), (C, D), (D, A)}
G = (V, E)
\end{lstlisting}

\(A\), \(B\), \(C\), and \(D\) are \textbf{vertices}. They may be drawn as points, but the graph itself does not require coordinate positions.

\subsection*{Example 3: Triangle mesh}

\begin{lstlisting}[basicstyle=\ttfamily\small, frame=single]
vertices = [
  { position: (0, 0, 0), normal: (0, 0, 1), uv: (0, 0) },
  { position: (1, 0, 0), normal: (0, 0, 1), uv: (1, 0) },
  { position: (0, 1, 0), normal: (0, 0, 1), uv: (0, 1) }
]
faces = [(0, 1, 2)]
\end{lstlisting}

Here each vertex contains a point-like \texttt{position}, but the vertex itself is a structured mesh element connected to a face.

\section{Recommended Terminology Standard}

Use the following convention in technical documentation:

\begin{enumerate}
  \item Use \textbf{point} for raw coordinate locations.
  \item Use \textbf{vertex} for graph elements and for coordinate-bearing elements that participate in geometry or topology.
  \item Use \textbf{node} when speaking in systems, networking, AST, tree, DOM, or infrastructure language, unless a graph-theoretic definition is needed.
  \item Avoid saying ``points and vertices are synonymous'' without a context qualifier.
\end{enumerate}

\callout{Best final wording}{In graph drawings, vertices are often depicted as points, but formally a graph vertex is an abstract object. In computational geometry, a point is a coordinate location, while a vertex is a point used structurally as part of a polygon, polyhedron, triangulation, graph embedding, or mesh.}

\section{References}

\begin{thebibliography}{9}

\bibitem{stanford_graph_basics}
Reza Zadeh, \emph{CME 305: Discrete Mathematics and Algorithms, Basic Definitions and Concepts in Graph Theory}, Stanford University. Defines a graph \(G(V,E)\) using a set of vertices and a set of edges. Available at: \url{https://stanford.edu/~rezab/discrete/Notes/2.pdf}

\bibitem{nagoya_graph_notes}
Richard Hepworth, \emph{Graph Theory}, Graduate School of Mathematics, Nagoya University, 2024. Gives a formal graph definition in terms of vertex and edge sets. Available at: \url{https://www.math.nagoya-u.ac.jp/~richard/teaching/s2024/Graph.pdf}

\bibitem{grinberg_graph_notes}
Darij Grinberg, \emph{An Introduction to Graph Theory}, 2026. Describes the vertex set notation \(V(G)\) for a graph. Available at: \url{https://www.cip.ifi.lmu.de/~grinberg/t/22s/graphs.pdf}

\bibitem{mount_comp_geom}
David M. Mount, \emph{CMSC 754: Computational Geometry Lecture Notes}, University of Maryland. Introduces computational geometry as algorithms for manipulating geometric objects. Available at: \url{https://graphics.stanford.edu/courses/cs268-09-winter/manuals/davidmount-754lects.pdf}

\bibitem{preparata_shamos_comp_geom}
Franco P. Preparata and Michael Ian Shamos, \emph{Computational Geometry: An Introduction}. Springer, 1985. Classic reference for computational geometry concepts and algorithms. Public scan available at: \url{https://www.rexresearch1.com/GeometryLibrary/ComputationalGeometryIntroductionPreparata.pdf}

\bibitem{mdn_webgl_vertexattribpointer}
MDN Web Docs, \emph{WebGLRenderingContext: vertexAttribPointer() method}. Describes binding buffer data to generic vertex attributes. Available at: \url{https://developer.mozilla.org/en-US/docs/Web/API/WebGLRenderingContext/vertexAttribPointer}

\bibitem{opengl_spec}
Mark Segal and Kurt Akeley, \emph{The OpenGL Graphics System: A Specification}. Khronos/OpenGL specification materials describe vertex-related graphics data including positions, normals, colors, and texture coordinates. Available at: \url{https://www.opengl.org/registry/doc/glspec20.20041022.pdf}

\end{thebibliography}

\end{document}
